% ---------------------------------------------------------------------
\section{Théorie}
\label{sec-theorie}
% ---------------------------------------------------------------------

Nous gardons la numérotation de l'énoncé : les équations (1) à (8) et les propriétés (1)
(translation) et (2) (filtrage) y renvoient directement. Nos propres équations sont
numérotées à partir de (9).

\subsection{Prétraitement et théorème de Shannon (Q1)}
\label{sec-shannon}

Le signal est ramené en mono (moyenne des canaux) puis rééchantillonné à
$\Fe = \qty{8}{\kilo\hertz}$.

\paragraph{Ce qu'impose Shannon.} Un signal échantillonné à $\Fe$ n'est représenté sans
ambiguïté que s'il ne contient aucune composante au-delà de $\Fe/2$ : une sinusoïde de
fréquence $f_0 \in\, ]\Fe/2, \Fe[$ donne exactement les mêmes échantillons qu'une
sinusoïde de fréquence $\Fe - f_0$. Choisir $\Fe = \qty{8}{\kilo\hertz}$, c'est donc décider
que seule la bande $[0, \qty{4}{\kilo\hertz}]$ servira, et s'obliger à éliminer tout ce qui
est au-dessus \emph{avant} de décimer.

\paragraph{Pourquoi la bande \qtyrange{0}{4}{\kilo\hertz} suffit.} L'identité d'un morceau
tient à sa mélodie, à son harmonie et à son rythme, c'est-à-dire aux fondamentales des
notes et à leurs premiers harmoniques. La fondamentale la plus aiguë du piano
(do\textsubscript{8}) vaut \qty{4186}{\hertz} ; celles des voix chantées et de la quasi-totalité
des instruments mélodiques restent sous \qty{2}{\kilo\hertz}. L'essentiel de l'énergie
musicale est concentré dans cette zone, et c'est là que se trouvent les maxima les plus
marqués du spectrogramme, donc les pics les plus stables. La bande téléphonique
(\qtyrange{300}{3400}{\hertz}) a été dimensionnée sur le même constat : elle suffit à
reconnaître une voix ou une mélodie.

\paragraph{Ce que l'on perd.} On perd la \og brillance \fg{} : harmoniques élevés,
cymbales, attaques, sifflantes. Trois raisons rendent cette perte acceptable :
\begin{enumerate}
  \item ce contenu est peu énergétique, donc il porte peu de pics, et des pics fragiles ;
  \item c'est celui que le canal détruit en premier : un petit haut-parleur coupe vers
        \qty{3.4}{\kilo\hertz}, l'absorption de l'air et des parois croît avec la fréquence,
        et le codeur MP3 supprime les hautes fréquences dès que le débit baisse. Une
        empreinte qui s'appuierait sur ces fréquences serait présente dans la référence et
        absente de la requête ;
  \item le coût diminue : le nombre d'échantillons est divisé par
        $\num{44100}/\num{8000} \approx \num{5.5}$, comme le nombre de canaux $K$ à
        résolution temporelle égale. Avec le facteur $\log_2 N$ de la FFT ($N$ passerait de
        1024 à environ 5600 pour garder $\Tw = \qty{128}{\milli\second}$), le coût de la
        TFCT baisse d'un facteur de l'ordre de 7.
\end{enumerate}

\paragraph{Pourquoi \texttt{resample} et pas une décimation naïve.} Garder un échantillon
sur $k$ sans filtrer replierait toute la bande \qtyrange{4}{22}{\kilo\hertz} dans
\qtyrange{0}{4}{\kilo\hertz}, en créant de faux pics. Le plus grave est que ce contenu
n'est pas le même dans la référence (fichier d'origine, aigus intacts) et dans la requête
(aigus coupés par le canal ou le MP3) : les faux pics différeraient des deux côtés et
détruiraient l'appariement sans cause visible. La fonction \texttt{resample} applique un
filtre passe-bas anti-repliement avant de décimer, ce qui supprime le problème.

\subsection{Transformée de Fourier à court terme et compromis temps-fréquence (Q2)}
\label{sec-tfct}

Le contenu musical évolue en permanence : une TFD calculée sur tout le morceau
mélangerait toutes les notes et ne dirait rien de leur succession. On calcule donc la
transformée de Fourier à court terme et le spectrogramme en décibels
\begin{equation}
  X[m,k] = \sum_{n=0}^{N-1} x[n+mH]\, w[n]\, e^{-j2\pi kn/\Nfft},
  \qquad
  \SdB[m,k] = 10\log_{10}\bigl(|X[m,k]|^2 + \varepsilon\bigr),
  \tag{2}
\end{equation}
avec une fenêtre de Hann périodique, $H = N/2$ et $\Nfft = N$. La trame $m$ et le canal $k$
correspondent à $t_m = mH/\Fe$ et $f_k = k\Fe/\Nfft$.

\paragraph{Résolutions.} La résolution temporelle est la durée de la fenêtre,
$\Tw = N/\Fe$. La résolution fréquentielle est la largeur à \qty{-3}{\decibel} du lobe
principal de la fenêtre, $\Delta f = \kw\,\Fe/N$, où $\kw$ ne dépend que de la forme de la
fenêtre. L'énoncé donne $\kw = \num{1.30}$, qui est la valeur de la fenêtre de
\emph{Hamming} ; pour la fenêtre de \emph{Hann} que nous utilisons, $\kw = \num{1.44}$
\cite{harris1978}. Le produit
\begin{equation}
  \Delta f \cdot \Tw = \kw
\end{equation}
ne dépend pas de $N$ : toute précision gagnée en fréquence est perdue en temps, et
inversement. Monter d'un demi-ton multiplie la fréquence par $2^{1/12}$ ; deux notes voisines
sont séparées si leur écart $f\,(2^{1/12}-1)$ dépasse $\Delta f$, donc au-dessus de
\begin{equation}
  f_{\min} = \frac{\Delta f}{2^{1/12}-1} \approx \num{16.8}\,\Delta f .
\end{equation}
Le nombre de trames par seconde vaut $\Fe/H = 2\Fe/N$. Le \cref{tab-resolutions} donne
ces grandeurs pour les quatre valeurs de $N$ demandées.

\begin{table}[h]
  \centering\small
  \caption{Résolutions de la TFCT à $\Fe = \qty{8}{\kilo\hertz}$, fenêtre de Hann
  ($\kw = \num{1.44}$), $H = N/2$. Dernière colonne : nombre $\Nt$ de valeurs entières de
  $\dm$ dans la zone cible $\Delta t \in [\num{0.1} ; 2]$\,s (\cref{sec-empreintes}).}
  \label{tab-resolutions}
  \begin{tabular}{@{}rrrrrr@{}}
    \toprule
    $N$ & $\Tw$ (ms) & $\Delta f$ (Hz) & trames/s & $f_{\min}$ (Hz) & $\Nt$ \\
    \midrule
    256  & 32  & 45          & \num{62.5}  & 757 & 119 \\
    512  & 64  & \num{22.5}  & \num{31.25} & 378 & 59 \\
    1024 & 128 & \num{11.25} & \num{15.6}  & 189 & 30 \\
    2048 & 256 & \num{5.6}   & \num{7.8}   & 95  & 15 \\
    \bottomrule
  \end{tabular}
\end{table}

\paragraph{Choix de $N = 1024$.} Quatre arguments vont dans le même sens.
\begin{enumerate}
  \item \emph{Fréquence.} Avec $\Delta f \approx \qty{11}{\hertz}$, deux demi-tons sont
        séparés dès \qty{190}{\hertz}, donc sur tout le registre mélodique. À $N = 512$, il
        faudrait monter à \qty{380}{\hertz} : les notes graves voisines se fondraient dans
        un même lobe et leurs pics deviendraient instables.
  \item \emph{Temps.} $\Tw = \qty{128}{\milli\second}$ reste plus court que la durée
        typique d'une note (\qtyrange{100}{500}{\milli\second}), si bien qu'une trame
        contient le plus souvent une seule configuration harmonique. À $N = 2048$, une
        trame de \qty{256}{\milli\second} mélangerait souvent deux notes successives.
  \item \emph{Spécificité des clés.} Le nombre de trames par seconde fixe le nombre $\Nt$
        de valeurs de $\dm$ possibles dans la zone cible, donc le nombre de clés $M$
        (\cref{sec-empreintes}). Passer à $N = 2048$ divise $\Nt$ et $M$ par deux.
  \item \emph{Robustesse au filtrage.} La propriété (2) exige que la réponse
        impulsionnelle soit courte devant la fenêtre (\cref{sec-proprietes}). Une fenêtre
        plus longue l'étend à des réponses plus longues : c'est le seul argument en faveur
        de $N$ plus grand, et il ne l'emporte pas sur les trois précédents.
\end{enumerate}
$N = 1024$ est le plus petit $N$ qui sépare les demi-tons sur tout le registre
mélodique ; tout $N$ plus grand dégrade la résolution temporelle et la spécificité sans
bénéfice musical.

\paragraph{Choix de $H = N/2$ et de $\Nfft = N$.} Avec $H = N/2$, les fenêtres de Hann
périodiques successives se recouvrent de moitié et leur somme est constante
($w[n] + w[n + N/2] = 1$) : chaque échantillon pèse autant dans l'analyse, et aucun
événement bref ne tombe dans un creux entre deux trames. Prendre $\Nfft > N$ (remplissage
par des zéros) interpolerait le spectre sans améliorer $\Delta f$, qui ne dépend que de $N$ :
on multiplierait les canaux, donc les clés possibles, sans information nouvelle, et les pics
passeraient plus facilement d'un canal à son voisin.

\subsection{Les deux propriétés qui fondent la méthode}
\label{sec-proprietes}

\paragraph{Propriété (1) : un décalage temporel translate le spectrogramme.}
Soit $y[n] = x[n - m_0H]$ avec $m_0$ entier. D'après (2),
\begin{equation}
\begin{split}
  X_y[m,k]
  &= \sum_{n=0}^{N-1} x[n + mH - m_0H]\, w[n]\, e^{-j2\pi kn/\Nfft} \\
  &= \sum_{n=0}^{N-1} x[n + (m-m_0)H]\, w[n]\, e^{-j2\pi kn/\Nfft}
  = X_x[m-m_0, k],
\end{split}
  \label{eq-translation}
\end{equation}
donc $S_y[m,k] = S_x[m-m_0,k]$ : le spectrogramme est translaté de $m_0$ trames. Le
résultat est exact, sans approximation, et vaut même pour la phase, parce que chaque
trame est analysée dans son propre repère temporel. C'est lui qui permet de retrouver la
position de l'extrait (\cref{sec-decision}).

\emph{Décalage quelconque.} Si $n_0 = m_0H + r$ avec $0 < r < H$, la trame $m$ de $y$
analyse un segment de $x$ qui commence $r$ échantillons avant la trame $m-m_0$ de sa
grille : elle tombe \emph{entre} deux trames de $x$, et~\eqref{eq-translation} ne
s'applique plus. Dans un passage stable, $|X|$ varie peu d'une position de fenêtre à
l'autre et les pics restent en place ; près d'une attaque ou d'une fin de note, le
maximum temporel d'un pic peut passer d'une trame à sa voisine. En rapportant chaque
trame de $y$ à la trame de $x$ la plus proche, l'écart est d'au plus $H/2 = N/4$
échantillons, soit un quart de fenêtre quel que soit $N$. Deux conséquences :
les votes du bon morceau se répartissent sur $\delta = m_0$ et $\delta = m_0 + 1$ au lieu
d'une seule case, ce qui abaisse le score ; et une paire dont un seul pic a migré voit son
$\dm$ changer d'une unité, donc sa clé est perdue.

\paragraph{Propriété (2) : un filtrage ajoute un biais qui ne dépend que de la fréquence.}
Soit $y = \alpha\,(h * x)$ avec $\alpha > 0$ et $h$ causal de longueur $L_h$. On prolonge
$w$ par zéro hors de $[0, N-1]$. En reportant $y[n+mH] = \alpha\sum_l h[l]\,x[n+mH-l]$
dans (2), puis en posant $n' = n - l$ :
\begin{align}
  X_y[m,k]
  &= \alpha \sum_{n} \sum_{l=0}^{L_h-1} h[l]\, x[n+mH-l]\, w[n]\, e^{-j2\pi kn/\Nfft} \nonumber\\
  &= \alpha \sum_{l=0}^{L_h-1} h[l]\, e^{-j2\pi kl/\Nfft}
     \sum_{n'} x[n'+mH]\, w[n'+l]\, e^{-j2\pi kn'/\Nfft} .
  \label{eq-filtrage-exact}
\end{align}
Si la fenêtre ne variait pas sur la durée de $h$, on pourrait remplacer $w[n'+l]$ par $w[n']$
et la somme intérieure deviendrait $X_x[m,k]$. On écrit donc $w[n'+l] = w[n'] + (w[n'+l] - w[n'])$ :
\begin{equation}
  X_y[m,k] = \alpha\,\hhat(f_k)\, X_x[m,k] + \varepsilon[m,k],
  \qquad
  \hhat(f_k) = \sum_{l=0}^{L_h-1} h[l]\, e^{-j2\pi f_k l/\Fe},
  \label{eq-filtrage}
\end{equation}
avec
\[
  \varepsilon[m,k] = \alpha \sum_{l} h[l]\, e^{-j2\pi kl/\Nfft}
  \sum_{n'} x[n'+mH]\,\bigl(w[n'+l] - w[n']\bigr)\, e^{-j2\pi kn'/\Nfft} .
\]
Le terme $\varepsilon$ vient de ce que la fenêtre ne voit pas exactement le même segment
de $x$ pour chaque retard $l$. Pour la fenêtre de Hann
$w[n] = \tfrac12\bigl(1 - \cos(2\pi n/N)\bigr)$, la pente est au plus $\pi/N$, et la fenêtre
est nulle à ses deux bords, donc continue une fois prolongée par zéro. Par conséquent
$|w[n'+l] - w[n']| \le \pi l/N$ pour tout $n'$, bords compris, et
\[
  |\varepsilon[m,k]| \;\le\; \alpha\,\frac{\pi}{N}
  \sum_{l} l\,|h[l]| \sum_{n'=-L_h}^{N-1} |x[n'+mH]| .
\]
Rapportée à $|X_x[m,k]|$, l'erreur est d'ordre $L_h/N$ : elle est négligeable si $L_h \ll N$.
(Avec une fenêtre rectangulaire, discontinue à ses bords, elle serait d'ordre 1 : la
continuité de la fenêtre de Hann est ici essentielle.) En prenant le module au carré puis le
logarithme, on obtient la relation (3) de l'énoncé :
\begin{equation}
  S_{y,\mathrm{dB}}[m,k] \approx S_{x,\mathrm{dB}}[m,k] + g(k),
  \qquad g(k) = 20\log_{10}\alpha + 20\log_{10}|\hhat(f_k)| .
  \tag{3}
\end{equation}
Le biais $g(k)$ ne dépend pas de la trame $m$ : dans le spectrogramme en décibels, le gain
et le filtrage de toute la chaîne se réduisent à un décalage vertical qui ne dépend que de
la fréquence. On ne garde que le module : la phase de $X_y$ contient un terme qui dépend
de $t_0$, inconnu, et la réverbération la rend aléatoire.

\paragraph{Vérification expérimentale (Q2).} Si (3) est vraie, $S_{y,\mathrm{dB}} - S_{x,\mathrm{dB}}$
tracé en fonction de $k$ donne la même courbe $g(k)$ pour toutes les trames, et cette
courbe suit $20\log_{10}|\hhat(f_k)|$. Pour un passe-bande de Butterworth
\qtyrange{300}{3400}{\hertz}, dont la réponse impulsionnelle décroît en quelques
millisecondes ($L_h \ll N = 1024$), les courbes doivent se superposer dans la bande
passante. Hors bande, $|X_y|$ tombe au niveau du bruit numérique et la différence devient
erratique : la relation ne vaut que là où le signal filtré reste au-dessus de ce plancher.
La \cref{fig-filtrage} le confirme : \aremplir{dispersion entre trames dans la bande
passante (en dB), écart à $20\log_{10}|\hhat|$, comportement hors bande}.

\begin{figure}[t]
  \centering
  \inclurefig[0.8\linewidth]{difference_spectres_passe_bande}{$S_{y,\mathrm{dB}} - S_{x,\mathrm{dB}}$ en fonction de k pour plusieurs trames, passe-bande 300-3400 Hz}
  \caption{Vérification de la propriété (2) : différence $S_{y,\mathrm{dB}} - S_{x,\mathrm{dB}}$
  en fonction de la fréquence pour \aremplir{5} trames d'un même morceau, après un filtre
  passe-bande \qtyrange{300}{3400}{\hertz}. En pointillés : $20\log_{10}|\hhat(f)|$.
  \aremplir{Ce qu'il faut y voir.}}
  \label{fig-filtrage}
\end{figure}

\subsection{Invariance des maxima locaux et condition (5)}
\label{sec-maxima}

La carte de constellation est l'ensemble des maxima locaux du spectrogramme (définition
(4), \cref{sec-pics}). Deux résultats justifient ce choix.

\paragraph{Invariance au gain.} Si $y = \alpha x$ avec $\alpha > 0$, alors
$X_y = \alpha X_x$ et $S_y = \alpha^2 S_x$. La fonction $s \mapsto \alpha^2 s$ est strictement
croissante, donc elle conserve l'ordre des valeurs : pour tout voisinage $V$,
\[
  S_y[m,k] = \max_{V(m,k)} S_y \iff S_x[m,k] = \max_{V(m,k)} S_x .
\]
L'ensemble des positions des maxima locaux est donc \emph{rigoureusement} identique pour
$x$ et $y$. On ne conserve que ces positions $(m,k)$, jamais les amplitudes, qui dépendent
de $\alpha$. Le seuil doit respecter la même invariance : un seuil local (médiane de $\SdB$
dans la bande, plus une marge) se décale de $20\log_{10}\alpha$ en même temps que les pics,
et un quota de pics par bande ne dépend que du classement des valeurs ; un seuil global
$\tau$ fixé en décibels ferait au contraire apparaître ou disparaître des pics avec le
volume. Seule la constante $\varepsilon$ de (2) échappe à l'invariance, mais elle
n'intervient que dans le silence numérique, où elle crée des plateaux de valeurs égales que
le test d'égalité prendrait pour des pics : le seuil local, strictement au-dessus de la
médiane, les élimine.

\paragraph{Invariance au filtrage : la condition (5).} Soit $(m,k_0)$ un pic de $S_x$ et
\[
  \mu = S_{x,\mathrm{dB}}[m,k_0] - \max_{(m',k') \in V^*} S_{x,\mathrm{dB}}[m',k'] \;>\; 0,
  \qquad V^* = V(m,k_0) \setminus \{(m,k_0)\},
\]
son avance en décibels sur son concurrent le plus fort dans le voisinage. Notons
\[
  \Delta g(V) = \max_{|k'-k_0| \le \kv} |g(k') - g(k_0)|
\]
la variation du biais sur la largeur du voisinage. D'après (3), pour tout
$(m',k') \in V^*$,
\[
  S_{y,\mathrm{dB}}[m,k_0] - S_{y,\mathrm{dB}}[m',k']
  = \underbrace{S_{x,\mathrm{dB}}[m,k_0] - S_{x,\mathrm{dB}}[m',k']}_{\ge\, \mu}
  + \underbrace{g(k_0) - g(k')}_{\ge\, -\Delta g(V)}
  \;\ge\; \mu - \Delta g(V) .
\]
Le pic reste donc le maximum de son voisinage dès que
\begin{equation}
  \Delta g(V) < \mu \quad \Longrightarrow \quad \text{le pic survit au filtrage.}
  \tag{5}
\end{equation}
La condition est suffisante, pas nécessaire : un pic peut survivre avec
$\Delta g(V) \ge \mu$ si son concurrent le plus proche en valeur n'est pas celui que le filtre
favorise. Elle est toujours remplie pour un gain seul ($g$ constant, $\Delta g = 0$). Pour un
filtre dont la réponse varie de $|g'|$ dB par canal, $\Delta g(V) \approx |g'|\,\kv$, qui ne
dépend que de la \emph{demi-largeur fréquentielle} $\kv$ du voisinage : la largeur
temporelle $\mv$ n'intervient pas, puisque $g$ ne dépend pas de $m$.

On en tire la règle de dimensionnement de l'énoncé. Un voisinage étroit en fréquence
donne un $\Delta g$ petit, donc des pics robustes au filtrage, mais il produit beaucoup de
pics, dont certains peu dominants ($\mu$ faible). Un voisinage large ne garde que des pics
très dominants, donc une empreinte compacte et spécifique, mais il exige
$|g'|\,\kv < \mu$ sur une plus grande largeur. Dans la bande passante du filtre
\qtyrange{300}{3400}{\hertz}, $g$ est presque plat et presque tous les pics survivent ; près
des fréquences de coupure, la pente atteint plusieurs décibels sur $\pm 5$ canaux, et seuls
les pics nettement dominants restent en place.

\subsection{Limites des deux propriétés}
\label{sec-limites}

\paragraph{Réverbération.} La propriété (2) suppose $L_h \ll N$. Une réverbération de temps
$T_{60}$ décroît de \qty{60}{\decibel} en $T_{60}$, donc de \qty{20}{\decibel} en $T_{60}/3$ ;
sa longueur utile est $L_h \approx T_{60}\Fe/3$, soit environ 800 échantillons pour
$T_{60} = \qty{0.3}{\second}$ (comparable à $N$ : cas limite) et \num{2700} pour
$T_{60} = \qty{1}{\second}$ (plus du double de $N$). La borne sur $\varepsilon$ n'est plus
petite, et la courbe $S_{y,\mathrm{dB}} - S_{x,\mathrm{dB}}$ dépend de la trame. Une meilleure
image est celle d'un \emph{bruit additif produit par le signal lui-même} : la queue
réverbérée d'une note se superpose aux trames suivantes. Elle comble les creux entre les
notes, efface les pics des notes brèves et laisse survivre ceux des notes tenues, plus
énergétiques que la traîne.

\paragraph{Décalage non entier.} La propriété (1) n'est exacte que pour un décalage multiple
de $H$ ; dans le cas réaliste, une partie des pics migre d'une trame
(\cref{sec-proprietes}).

\paragraph{Tempo et hauteur.} Ni (1) ni (2) ne couvrent un changement de tempo, qui
modifie les écarts $\dm$, ni un changement de hauteur, qui multiplie toutes les fréquences
par un même facteur et change donc les canaux $k$ et les écarts $\dk$ : la méthode n'y est
pas invariante.
